\documentclass[11pt]{article}

\usepackage[margin=1.05in]{geometry}
\usepackage[T1]{fontenc}
\usepackage{lmodern}
\usepackage{microtype}
\usepackage{amsmath,amssymb,amsthm,mathtools}
\usepackage{enumitem}
\usepackage{booktabs}
\usepackage[hidelinks]{hyperref}

\setlist{nosep,leftmargin=*}
\setlength{\parindent}{1.25em}
\setlength{\parskip}{0.15em}
\numberwithin{equation}{section}

\newtheorem{theorem}{Theorem}[section]
\newtheorem{proposition}[theorem]{Proposition}
\newtheorem{lemma}[theorem]{Lemma}
\newtheorem{corollary}[theorem]{Corollary}
\theoremstyle{definition}
\newtheorem{definition}[theorem]{Definition}
\newtheorem{example}[theorem]{Example}
\theoremstyle{remark}
\newtheorem{remark}[theorem]{Remark}

\newcommand{\C}{\mathbb C}
\newcommand{\Q}{\mathbb Q}
\newcommand{\R}{\mathbb R}
\newcommand{\Z}{\mathbb Z}
\newcommand{\N}{\mathbb N}
\newcommand{\HH}{\mathbb H}
\newcommand{\Span}{\operatorname{span}}
\newcommand{\Tr}{\operatorname{Tr}}
\newcommand{\dtrdeg}{\operatorname{dtrdeg}}
\newcommand{\Tor}{\operatorname{Tor}}
\newcommand{\dd}{\mathrm d}
\newcommand{\T}{\mathsf T}
\newcommand{\Deltaop}{\mathop{\Delta}}

\title{\bfseries Differential Algebraicity of Multiple Elliptic Gamma
Functions:\\Nome-Rank Classification and Translation-Coset Matroids}
\author{Michael Haufschild}
\date{July 22, 2026}

\begin{document}
\maketitle

\begin{abstract}
Let $I$ be a nonempty finite set, let $\tau=(\tau_i)_{i\in I}$ lie in the
complex upper half-plane, and let $G_I(z;\tau)$ be Nishizawa's multiple
elliptic gamma function.  Define
\[
 d_{\Q}(\tau)=
 \dim_{\Q}\Span_{\Q}\bigl(\{1\}\cup\{\tau_i:i\in I\}\bigr).
\]
We prove that $G_I$ is differentially algebraic over $\C(z)$ if and only if
$d_{\Q}(\tau)=2$, and is hypertranscendental if and only if
$d_{\Q}(\tau)\geq3$.  The same dichotomy holds for its logarithmic
derivative.  Thus every finite upper-half-plane period tuple is classified,
including partially dependent non-rational tuples.

The invariant has an intrinsic multiplicative form:
$d_{\Q}(\tau)-1$ is the torsion-free rank of the group generated by the
nomes $e^{2\pi i\tau_i}$.  For the rank-three elliptic-gamma seed we prove
more.  Put $P=\Z+\Z\sigma+\Z\tau$, let $K$ be the compositum of the two
elliptic function fields, and set
$D_h=\Deltaop_h(\partial\Gamma/\Gamma)$ and
$R_h=\Gamma(z+h)/\Gamma(z)$.  For every finite shift family and for
$B=K$ or $B=\C(z)$, the differential transcendence degree of either the
$D_h$-family or the $R_h$-family over $B$ is exactly the number of distinct
nonzero classes in $\C/P$.  Thus the entire shift packet is the pointed
partition matroid of period classes, and $\Gamma$ is hypertranscendental even
over $K$.

The positive phase follows from an all-coordinate multiplication identity
and Kato's rational-period theorem.  The negative phase uses joint Ostrowski
descent, product-curve specialization, and an orbital pole obstruction.
Nishizawa's face equations propagate the seed to every higher order.  Together
these results close Kato's problem for non-rational periods and determine all
finite translation dependencies of the seed.
\end{abstract}

\noindent\textbf{2020 Mathematics Subject Classification.}
Primary 33E30; Secondary 12H05, 39A13.

\noindent\textbf{Keywords.}
Multiple elliptic gamma function; differential algebraicity;
hypertranscendence; elliptic translation; nome rank; differential matroid;
Ostrowski descent.

\section{Introduction}

For periods $\tau_i$ in the upper half-plane
$\HH=\{w\in\C:\operatorname{Im}w>0\}$, put
\[
 q_i=e^{2\pi i\tau_i},\qquad x=e^{2\pi iz},\qquad
 Q_I=\prod_{i\in I}q_i.
\]
For a finite family $q=(q_i)_{i\in I}$, write
\[
 (x;q)_\infty=
 \prod_{n\in\N^I}\left(1-x\prod_{i\in I}q_i^{n_i}\right).
\]
The standard Nishizawa function is
\begin{equation}
 G_I(z;\tau)=
 (Q_I/x;q)_\infty\,
 (x;q)_\infty^{\,(-1)^{|I|-1}}.
 \label{eq:defG}
\end{equation}
When $I=\{0,\ldots,r\}$, this is Nishizawa's notation
$G_r(z;(\tau_0,\ldots,\tau_r))$.
The products converge locally uniformly because $|q_i|<1$.  Formula
\eqref{eq:defG} therefore defines a meromorphic function of $z$ on $\C$.
For $|I|=1$ it is a theta function; for $|I|=2$ it is the elliptic gamma
function.  These functions occur in elliptic hypergeometric analysis,
integrable models, supersymmetric partition functions, and modular and
arithmetic constructions; see
\cite{FelderVarchenko2000,Narukawa2004,Nishizawa2001}.

We equip $\C(z)$ and fields of meromorphic functions with
$\partial=\dd/\dd z$.  A meromorphic function is
\emph{differentially algebraic} over $\C(z)$ if a nonzero polynomial over
$\C(z)$ vanishes on a finite jet
$(f,\partial f,\ldots,\partial^m f)$.  It is
\emph{hypertranscendental} if it is not differentially algebraic.

Kato proved differential algebraicity when the period vector is rational:
there is an $\omega\in\HH$ such that every $\tau_i$ is a positive rational
multiple of $\omega$ \cite[Theorem~1.1]{Kato2021}.  He then stated that the
non-rational case was unknown and anticipated possible exceptional cases
\cite[Remark~1.2]{Kato2021}.  Literal rational-ray membership is not the
sharp invariant.  Integer offsets disappear under exponentiation, and
finite root-of-unity twists disappear after a multiplication identity.  The
correct invariant is the augmented rational rank
\begin{equation}
 d_{\Q}(\tau)=
 \dim_{\Q}\Span_{\Q}\bigl(\{1\}\cup\{\tau_i:i\in I\}\bigr).
 \label{eq:rank}
\end{equation}
The unit period appears because \eqref{eq:defG} is $1$-periodic.  Since $I$
is nonempty and every $\tau_i$ is nonreal, $d_{\Q}(\tau)\geq2$.

Our result gives a complete answer for the standard finite Nishizawa
hierarchy.  Rank two is exactly the differential-algebraic phase.  Every
higher rank is hypertranscendental.  The difficult seed is the
two-period elliptic gamma with rationally independent periods
$1,\sigma,\tau$.  We prove this seed without assuming a
Picard--Vessiot existence or constant-extension theorem.  Instead, a short
K\"ahler-differential argument gives an exact finite-family criterion for
simultaneous telescoping; a trace argument descends its primitive to the base
field.  The finite-family form then identifies every differential dependence
among translates of the seed, not only its individual differential nature.

There are two main results.  Theorem~\ref{thm:main} gives the universal
rational-rank, equivalently nome-rank, classification for the whole finite
hierarchy.  Theorem~\ref{thm:coset-matroid} gives the exact translation-coset
differential matroid of the rank-three seed.  Multiplication formulas and
single-orbit summability obstructions are established mechanisms
\cite{FelderVarchenkoMultiplication,DreyfusHardouinRoquesSinger2018,
ArrecheBabbitt2025}; the new content is their synthesis with simultaneous
descent into these two complete classifications.

\section{Main theorem}

Put
\[
 \Psi_I(z;\tau)=\frac{\partial G_I(z;\tau)}{G_I(z;\tau)}.
\]

\begin{theorem}[Exact rational-rank dichotomy]
\label{thm:main}
Let $I$ be a nonempty finite set and let $\tau:I\to\HH$.  Then
\begin{align*}
 G_I(\,\cdot\,;\tau)\text{ is differentially algebraic over }\C(z)
 &\Longleftrightarrow d_{\Q}(\tau)=2,\\
 G_I(\,\cdot\,;\tau)\text{ is hypertranscendental over }\C(z)
 &\Longleftrightarrow d_{\Q}(\tau)\geq3.
\end{align*}
Both equivalences remain true with $G_I$ replaced by $\Psi_I$.
\end{theorem}

The two alternatives in Theorem~\ref{thm:main} are exhaustive because
$d_{\Q}(\tau)\geq2$.  No independence is required among all periods in the
hypertranscendental phase: two independent period directions suffice.

\begin{corollary}[Nome-rank law]
\label{cor:nome-rank}
Let $q_i=e^{2\pi i\tau_i}$ and let $\mathcal Q_\tau\subset\C^\times$ be the
multiplicative subgroup generated by the $q_i$.  Then
\begin{equation}
 d_{\Q}(\tau)-1=
 \operatorname{rank}_{\Z}(\mathcal Q_\tau/\Tor(\mathcal Q_\tau)).
 \label{eq:nome-rank}
\end{equation}
Consequently $G_I$ and $\Psi_I$ are differentially algebraic exactly when
the nomes generate a group of torsion-free rank one, and they are
hypertranscendental exactly when that rank is at least two.
\end{corollary}

\begin{proof}
Modulo the line $\Q\cdot1$, the rational span of the periods has dimension
$d_{\Q}(\tau)-1$.  A multiplicative relation
$\prod_iq_i^{n_i}$ is torsion precisely when
$m\sum_i n_i\tau_i\in\Z$ for some $m\geq1$, equivalently when
$\sum_i n_i\tau_i\in\Q$.  Tensoring the finitely generated group
$\mathcal Q_\tau/\Tor(\mathcal Q_\tau)$ with $\Q$ therefore identifies it
with
\[
 \Span_{\Q}\{1,\tau_i:i\in I\}/\Q\cdot1.
\]
This proves \eqref{eq:nome-rank}; Theorem~\ref{thm:main} gives the final
assertions.
\end{proof}

\begin{corollary}[Kato's non-rational-period problem]
\label{cor:kato}
A period vector outside Kato's rational-ray locus is differentially
algebraic precisely when its augmented rational rank is two.  It is
hypertranscendental precisely when its augmented rational rank is at least
three.
\end{corollary}

\begin{example}[The exceptional non-rational phase]
\label{ex:rank2}
Fix $\omega\in\HH$.  Every tuple
\[
 \tau_i=a_i+b_i\omega,\qquad
 a_i,b_i\in\Q,\quad b_i>0,
\]
has augmented rank two.  For every $\omega\in\HH$, the tuple
$(\omega,\omega+\tfrac12,2\omega+\tfrac13)$ is not contained in a rational
ray, but Theorem~\ref{thm:main} makes its multiple elliptic gamma function
differentially algebraic.
\end{example}

\begin{example}[Dependent hypertranscendental phase]
If distinct labels $q,r$ make $1,\tau_q,\tau_r$ rationally independent,
then arbitrary dependencies among the remaining periods are harmless.
Repeated periods, rational multiples of $\tau_q$, and periods in the
rational span of earlier entries do not restore differential algebraicity.
\end{example}

\begin{remark}[Scope]
Theorem~\ref{thm:main} concerns the standard one-variable Nishizawa product
\eqref{eq:defG} for arbitrary nonempty finite upper-half-plane period
families.  Cone-indexed and geometric-family generalizations
\cite{MorainI,MorainII,Winding2016} are different functions, not omitted
parameter cases of this theorem.
\end{remark}

\section{Differential-algebraic preliminaries}

For a differential field $B$ inside a differential overfield and a set $S$,
write $B\langle S\rangle_\partial$ for the smallest differential subfield
containing $B\cup S$.  We use the following standard closure facts in their
field-theoretic form.

\begin{lemma}[Closure, logarithmic derivatives, and base change]
\label{lem:da-preliminaries}
Let $B\subseteq L\subseteq\mathcal M$ be differential fields of
characteristic zero.
\begin{enumerate}[label=\textup{(\roman*)}]
\item Elements differentially algebraic over $B$ form a differential
subfield of $\mathcal M$, and differential algebraicity is transitive in
towers.
\item If $0\ne f\in\mathcal M$, then $f$ is differentially algebraic over
$B$ if and only if $f'/f$ is.
\item If $L/B$ is differentially algebraic, then for every finite
$S\subset\mathcal M$,
\[
 \dtrdeg_B B\langle S\rangle_\partial
 =\dtrdeg_L L\langle S\rangle_\partial.
\]
\end{enumerate}
\end{lemma}

\begin{proof}
Part (i) is the differential-field closure theorem and the tower theorem for
differential transcendence degree.  For (ii), one direction follows from
field closure.  Conversely, if $a=f'/f$ is differentially algebraic over
$B$, then $B\langle a\rangle_\partial/B$ is differentially algebraic and
$f$ satisfies $Y'-aY=0$ over that field; transitivity applies.  For (iii),
apply the tower formula twice inside $L\langle S\rangle_\partial$.  Both
$L/B$ and $L\langle S\rangle_\partial/B\langle S\rangle_\partial$ have
differential transcendence degree zero, so the remaining two terms agree.
\end{proof}

For $c\in\C$, translation $\T_c f(z)=f(z+c)$ commutes with $\partial$.
It need not fix a coefficient field pointwise: throughout the elliptic
argument it acts \emph{semilinearly}, as an automorphism stabilizing that
field.  Such an automorphism preserves differential algebraicity and all
differential dependence relations.  We will also use that $B(z)/B$ is
differentially algebraic because $z'-1=0$.

\section{The complete rank-two phase}

\begin{lemma}[Affine rational normal form]
\label{lem:normal}
If $d_{\Q}(\tau)=2$ and $\omega=\tau_{i_0}$ for one $i_0\in I$, then
\begin{equation}
 \tau_i=a_i+b_i\omega,\qquad
 a_i,b_i\in\Q,\quad b_i>0
 \label{eq:affine}
\end{equation}
for every $i\in I$.
\end{lemma}

\begin{proof}
The pair $\{1,\omega\}$ is rationally independent because
$\operatorname{Im}\omega>0$.  It is therefore a basis of the augmented
period span.  Taking imaginary parts in the resulting rational expansion
gives
$\operatorname{Im}\tau_i=b_i\operatorname{Im}\omega$, hence $b_i>0$.
\end{proof}

Choose a positive integer $M$ clearing all denominators in
\eqref{eq:affine}.  Then
\begin{equation}
 M\tau_i=A_i+B_i\omega,\qquad
 A_i\in\Z,\quad B_i\in\Z_{>0},
 \label{eq:integer-normal}
\end{equation}
so the corresponding nomes satisfy
\begin{equation}
 e^{2\pi iM\tau_i}=e^{2\pi iB_i\omega}.
 \label{eq:nome-normal}
\end{equation}

\begin{lemma}[All-coordinate multiplication identity]
\label{lem:multiplication}
For every positive integer $M$,
\begin{equation}
 G_I(z;\tau)=
 \prod_{k\in\prod_{i\in I}\{0,\ldots,M-1\}}
 G_I\left(z+\sum_{i\in I}k_i\tau_i;M\tau\right).
 \label{eq:multiplication}
\end{equation}
\end{lemma}

This residue-class factorization is the simultaneous-coordinate form needed
here of the established elliptic-gamma multiplication mechanism; compare
Felder--Varchenko \cite{FelderVarchenkoMultiplication}.  We include the proof
to fix conventions and because the arbitrary finite coordinate set is used
in the classification.

\begin{proof}
Write $p_i=q_i^M$ and
$x_k=x\prod_iq_i^{k_i}$.  In every factor
$(x_k;p)_\infty$, the exponent of $q_i$ is
$M\ell_i+k_i$.  As $\ell_i\geq0$ and $0\leq k_i<M$, these exponents
partition $\N$ uniquely.  Hence the product of all $x$-side
$q$-shifted factorials on the right of \eqref{eq:multiplication} is
$(x;q)_\infty$.

For the $Q_I/x$-side factorial, the exponent of $q_i$ is
\[
 M(\ell_i+1)-k_i.
\]
As $(\ell_i,k_i)$ varies, this partitions the positive integers.  Equivalently,
write
\[
 M(\ell_i+1)-k_i=M\ell_i+(M-1-k_i)+1.
\]
Thus the product of all $Q_I/x$-side factorials is
$(Q_I/x;q)_\infty$.  The parity exponent
$(-1)^{|I|-1}$ is unchanged.  Local uniform absolute convergence permits
these finite residue decompositions, proving \eqref{eq:multiplication} as a
meromorphic identity.
\end{proof}

\begin{proposition}[Rank two is differentially algebraic]
\label{prop:rank2}
If $d_{\Q}(\tau)=2$, then $G_I$ and $\Psi_I$ are differentially algebraic
over $\C(z)$.
\end{proposition}

\begin{proof}
Apply Lemma~\ref{lem:multiplication} with the denominator-clearing integer
$M$.  By \eqref{eq:nome-normal} and the product definition, the function
with period tuple $M\tau$ is identical to the function with tuple
$(B_i\omega)_{i\in I}$.  The latter is differentially algebraic by
Kato's theorem \cite[Theorem~1.1]{Kato2021}.

For each $c\in\C$, translation $z\mapsto z+c$ is a differential
automorphism of $\C(z)$ and preserves differential algebraicity of
meromorphic functions.  Differentially algebraic elements form a field.
The finite product in \eqref{eq:multiplication} is therefore
differentially algebraic.  Closure under differentiation and division by a
nonzero element gives the same conclusion for
$\Psi_I=\partial G_I/G_I$.
\end{proof}

\begin{remark}[Constructive content of the rank-two proof]
\label{rem:effective}
The reduction is explicit: from rational coefficients $a_i,b_i$ it computes
$M,A_i,B_i$ and expresses $G_I$ as exactly $M^{|I|}$ translates of the
Kato-rational function with periods $(B_i\omega)_i$.  Over any specified
computable differential coefficient field, annihilators for those factors
can therefore be combined by standard differential elimination into an
annihilator for $G_I$.  An unconditional algorithm for arbitrary
$\omega\in\HH$ is not a well-posed strengthening without an effective
representation of $\omega$ and an effective version of the Kato input.  The
present theorem is universal in arbitrary complex periods and leaves no
parameter case unclassified; its constructive boundary is precisely the
positive theorem it invokes.
\end{remark}

\section{An exact joint Ostrowski criterion}

We isolate the algebraic mechanism needed for the elliptic-gamma seed in a
finite-family form.  No commutation law is required among the automorphisms.
The one-automorphism parameterized Picard--Vessiot criterion is classical;
see Hardouin--Singer \cite[Proposition~3.1]{HardouinSinger2008}.  The useful
form here is an elementary simultaneous criterion in which one trace produces
one common base-field primitive in every direction.

\begin{theorem}[Joint Ostrowski criterion]
\label{lem:ostrowski}
Let $C\subseteq k\subseteq F$ be fields of characteristic zero.  Suppose
$F$ is a differential field, $k$ is stable under $\partial$, and $C$
consists of differential constants.  Let $(\rho_j)_{j\in J}$ be any family
of automorphisms of $F$ which stabilize $k$, fix $C$, and commute with
$\partial$.  Assume
\[
 F^{(\rho_j)_{j\in J}}=C.
\]
For $1\leq r\leq m$, let $y_r\in F$ satisfy
\[
 \rho_j(y_r)-y_r=a_{j,r}\in k\qquad(j\in J).
\]
Then $y_1,\ldots,y_m$ are differentially algebraically dependent over $k$
if and only if there are operators $L_1,\ldots,L_m\in C[\partial]$, not all
zero, and an element $g\in k$ such that
\begin{equation}
 \rho_j(g)-g=\sum_{r=1}^m L_r(a_{j,r})
 \qquad(j\in J).
 \label{eq:joint-telescope}
\end{equation}
Whenever this holds, there is a $c\in C$ such that
\begin{equation}
 \sum_{r=1}^m L_r(y_r)=g+c.
 \label{eq:joint-linear-relation}
\end{equation}
\end{theorem}

\begin{proof}
Suppose first that the family is differentially algebraically dependent.
There is an $N$ such that the finite jet family
$x_{r,n}=\partial^ny_r$ $(1\leq r\leq m,\ 0\leq n\leq N)$ is algebraically
dependent over $k$.  Its K\"ahler differentials in $\Omega_{F/k}$ are
therefore linearly dependent.  Let
\[
 W=\left\{(u_{r,n})\in F^{m(N+1)}:
          \sum_{r,n}u_{r,n}\,\dd x_{r,n}=0\right\}.
\]
This is a nonzero $F$-subspace.  Since
\[
 \rho_j(x_{r,n})=x_{r,n}+\partial^na_{j,r}
\]
and $\partial^na_{j,r}\in k$, applying $\rho_j$ to a relation in $W$ gives
another relation in $W$.  Explicitly, because $\rho_j(k)=k$, it induces the
semilinear automorphism
\[
 (\rho_j)_*:\Omega_{F/k}\longrightarrow\Omega_{F/k},\qquad
 a\,\dd x\longmapsto\rho_j(a)\,\dd(\rho_jx),
\]
and the displayed increment has zero relative differential.  Thus $W$ is
stable under the coordinatewise semilinear action.

Choose a nonzero vector $u\in W$ with minimal support and divide by one
nonzero coordinate so that this coordinate equals $1$.  For each $j$, the
vector $\rho_j(u)-u$ belongs to $W$ and has strictly smaller support.
Minimality forces $\rho_j(u)=u$ for every $j\in J$.  Every coordinate of
$u$ consequently belongs to the joint fixed field $C$.  We obtain
coefficients $c_{r,n}\in C$, not all zero, such that, for
$L_r=\sum_n c_{r,n}\partial^n$,
\[
 \dd_{F/k}\!\left(\sum_rL_r(y_r)\right)=0
 \quad\text{in }\Omega_{F/k}.
\]
In characteristic zero, the kernel of the universal derivation
$\dd_{F/k}:F\to\Omega_{F/k}$ is the relative algebraic closure of $k$ in
$F$.
Indeed, if an element of that kernel were transcendental over $k$, extend
it to a transcendence basis of $F/k$, send it to $1$ and the other basis
elements to $0$, and extend this derivation uniquely across the remaining
algebraic extension.  This contradicts membership in the kernel.
Thus $u_0:=\sum_rL_r(y_r)$ is algebraic over $k$.

Put $E=k(u_0)$ and $e=[E:k]$.  Because
\[
 \rho_j(u_0)-u_0=\sum_rL_r(a_{j,r})\in k,
\]
each $\rho_j$ maps $E$ onto itself: the displayed formula gives one
inclusion and its rearrangement gives the reverse inclusion.  Define
\[
 g=\frac1e\Tr_{E/k}(u_0)\in k.
\]
Trace is compatible with the semilinear action of $\rho_j$.  Indeed,
applying $\rho_j$ to a multiplication matrix for $u_0$ gives a multiplication
matrix for $\rho_j(u_0)$, so
$\Tr_{E/k}(\rho_j u_0)=\rho_j(\Tr_{E/k}u_0)$.  The trace of an element of
$k$ is $e$ times that element.  Therefore
\[
 \rho_j(g)-g
 =\frac1e\Tr_{E/k}\bigl(\rho_j(u_0)-u_0\bigr)
 =\sum_rL_r(a_{j,r}).
\]
It follows that $u_0-g$ is fixed by every $\rho_j$, hence belongs to $C$.
This proves \eqref{eq:joint-telescope} and
\eqref{eq:joint-linear-relation}.

Conversely, suppose \eqref{eq:joint-telescope} holds.  Then
$\sum_rL_r(y_r)-g$ is fixed by every $\rho_j$, so it lies in $C$.  Since
the operators are not all zero, this is a nontrivial differential-polynomial
relation over $k$.  The family $(y_r)$ is differentially algebraically
dependent.
\end{proof}

\begin{remark}
Theorem~\ref{lem:ostrowski} requires neither a differential closure of $C$
nor a Picard--Vessiot existence theorem.  One field trace produces the same
base element $g$ in every direction.  The proof also shows that the index
set $J$ may be infinite: every dependence still yields finitely many
operators because it involves a finite jet family.
\end{remark}

\section{The rank-three translation-coset matroid}

Let $\sigma,\tau\in\HH$, set
\[
 \Gamma(z)=G_{\{0,1\}}(z;\sigma,\tau),\qquad
 \theta_\omega(z)=G_{\{0\}}(z;\omega),
\]
and put
\[
 \Psi=\frac{\partial\Gamma}{\Gamma},\qquad
 \lambda_\omega=\frac{\partial\theta_\omega}{\theta_\omega},\qquad
 \Deltaop_hf=\T_hf-f.
\]
For $p=e^{2\pi i\sigma}$ and $q=e^{2\pi i\tau}$, direct residue
separation in \eqref{eq:defG} gives
\[
 \frac{\Gamma(z+\sigma)}{\Gamma(z)}
  =(q/x;q)_\infty(x;q)_\infty=\theta_\tau(z),
\]
and the same calculation with $p$ and $q$ interchanged gives the
$\tau$-shift.  Thus
\begin{equation}
 \Gamma(z+1)=\Gamma(z),\quad
 \Gamma(z+\sigma)=\theta_\tau(z)\Gamma(z),\quad
 \Gamma(z+\tau)=\theta_\sigma(z)\Gamma(z),
 \label{eq:gamma-equations}
\end{equation}
and
\begin{equation}
 \Deltaop_1\Psi=0,\qquad
 \Deltaop_\sigma\Psi=\lambda_\tau,\qquad
 \Deltaop_\tau\Psi=\lambda_\sigma.
 \label{eq:psi-equations}
\end{equation}

Write $\Lambda_\omega=\Z+\Z\omega$ and let $E_\omega$ be the field of
$\Lambda_\omega$-elliptic meromorphic functions.  The theta function has
one simple zero modulo $\Lambda_\omega$.  Its lattice shifts multiply it by
nonzero exponential factors, so shifts of $\lambda_\omega$ differ from it
by constants.  Hence $\lambda_\omega$ has one simple pole, of residue one,
modulo $\Lambda_\omega$, while
\begin{equation}
 \lambda_\omega'\in E_\omega,\qquad
 \Deltaop_h\lambda_\omega\in E_\omega.
 \label{eq:lambda-elliptic}
\end{equation}

Assume from now on that $1,\sigma,\tau$ are $\Q$-linearly independent and
define
\[
 P=\Z+\Z\sigma+\Z\tau,\qquad K=E_\tau E_\sigma,
\]
\begin{equation}
 D_h=\Deltaop_h\Psi,\qquad
 R_h=\frac{\Gamma(z+h)}{\Gamma(z)}.
 \label{eq:defects-quotients}
\end{equation}

We record all geometric inputs used below.

\begin{lemma}[Nondiscrete periods]
\label{lem:nondiscrete}
The subgroup $P\subset\C\simeq\R^2$ is nondiscrete.  A meromorphic function
whose period group contains $P$ is constant.
\end{lemma}

\begin{proof}
A discrete subgroup of $\R^2$ has rank at most two, whereas rational
independence makes $P$ free of rank three.  Thus $P$ contains nonzero
elements tending to zero.  A meromorphic function with such periods has no
pole, since translates of a pole would accumulate.  It is entire, and
difference quotients along the vanishing period sequence force its
derivative to vanish.
\end{proof}

\begin{lemma}[Elliptic product field and differential bases]
\label{lem:elliptic-fields}
The fields $E_\tau$ and $E_\sigma$ are algebraically disjoint over $\C$.
Consequently $K$ is the function field of
$(\C/\Lambda_\tau)\times(\C/\Lambda_\sigma)$.  Moreover $K/\C$ is
differentially algebraic, so $K(z)$ is differentially algebraic over both
$K$ and $\C(z)$.
\end{lemma}

\begin{proof}
If algebraic disjointness failed, regularity of the two one-variable
function fields would give algebraically dependent nonconstant functions
$f\in E_\tau$ and $v\in E_\sigma$.  Translation by $\tau$ fixes $f$, so all
$v(z+n\tau)$ satisfy one polynomial over $\C(f)$.  A polynomial has only
finitely many roots in a field; hence $v(z+N\tau)=v(z)$ for some nonzero
integer $N$.  The period group of $v$ then contains
$\Z+\Z\sigma+\Z N\tau$, a rank-three subgroup.  The proof of
Lemma~\ref{lem:nondiscrete} makes $v$ constant, a contradiction.  Therefore
$E_\tau\otimes_\C E_\sigma$ is a domain, and its fraction field is the
function field of the product curve.

For each lattice, the elliptic function field is generated by its
Weierstrass function and derivative, which satisfy the Weierstrass algebraic
differential equation.  Thus each $E_\omega/\C$, and hence $K/\C$, is
differentially algebraic.  Finally $z'-1=0$ gives the two assertions about
$K(z)$.
\end{proof}

\begin{lemma}[Generic-fiber specialization]
\label{lem:specialization}
Let $X,Y$ be irreducible curves, let $a$ be an automorphism of $X$, and let
$r\in\C(X\times Y)$.  If
\[
 (a\times1)^*r-r=s\in\C(X),
\]
then there is a closed point $y\in Y$ at which $r$ restricts to
$r_y\in\C(X)$ and $(a^*-1)r_y=s$.
\end{lemma}

\begin{proof}
Write $r$ as a quotient on affine opens of $X$ and $Y$.  Excluding the
proper closed subset of $Y$ on which its denominator specializes to zero
leaves a nonempty open set of admissible fibers.  Restriction to any closed
point in that open set is a field homomorphism on the coefficients involved,
commutes with the automorphism acting on $X$, and fixes $s$.  It therefore
preserves the displayed identity.
\end{proof}

\begin{lemma}[One-orbit boundary obstruction]
\label{lem:pole}
Let $X$ be an elliptic curve, let $\T_a$ be translation by a nontorsion
point, and let $u\in\C(X)$.  On every $\T_a$-orbit, the polar support of
$(\T_a-1)u$ is either empty or contains at least two points.
\end{lemma}

\begin{proof}
The polar divisor of $u$ is finite.  On a fixed orbit, index its finite
nonempty polar support by integers.  At the point immediately before the
least index, exactly the translated summand is singular.  At the greatest
index, exactly the untranslated summand is singular.  The points are
distinct because $a$ is nontorsion.  Cancellation is impossible at both
boundaries.
\end{proof}

This is the elementary support consequence of elliptic orbital residues;
compare Dreyfus--Hardouin--Roques--Singer
\cite[Appendix~B, Proposition~B.8]{DreyfusHardouinRoquesSinger2018} and the
complete orbital/panorbital summability criterion of Arreche--Babbitt
\cite[Theorems~1.1 and~2.6]{ArrecheBabbitt2025}.

\begin{lemma}[Highest-pole persistence]
\label{lem:highest-pole}
If $f$ has a simple pole at $x$ and
$0\ne L=\sum_{j=0}^N c_j\partial^j\in\C[\partial]$ with $c_N\ne0$, then
$L(f)$ has a nonzero pole of order $N+1$ at $x$.  It has no poles away from
those of $f$.
\end{lemma}

\begin{proof}
If the principal part of $f$ is $a(z-x)^{-1}$ with $a\ne0$, then the term
$c_N\partial^Nf$ has leading coefficient
$c_N(-1)^NN!a$ in degree $-(N+1)$.  Lower derivatives cannot cancel it.
Differentiation creates no new pole locations.
\end{proof}

\begin{theorem}[Translation-coset differential-rank law]
\label{thm:coset-matroid}
For shifts $h_1,\ldots,h_m\in\C$, let $\nu$ be the number of distinct
nonzero classes among $h_1+P,\ldots,h_m+P$ in $\C/P$.  For each
$B\in\{K,\C(z)\}$,
\begin{equation}
 \dtrdeg_B B\langle D_{h_1},\ldots,D_{h_m}\rangle_\partial
 =\dtrdeg_B B\langle R_{h_1},\ldots,R_{h_m}\rangle_\partial
 =\nu.
 \label{eq:coset-rank}
\end{equation}
Equivalently, the zero period class is a loop, each nonzero period class is
a rank-one parallel class, and distinct nonzero classes are freely
differentially independent.
\end{theorem}

\begin{proof}
First work over $K$.  The translation cocycle identities are
\begin{equation}
 D_{u+v}=\T_uD_v+D_u,\qquad D_{-u}=-\T_{-u}D_u.
 \label{eq:defect-cocycle}
\end{equation}
By \eqref{eq:psi-equations}, $D_1=0$, $D_\sigma=\lambda_\tau$, and
$D_\tau=\lambda_\sigma$.  Equation \eqref{eq:lambda-elliptic} makes the
last two functions differentially algebraic over $K$.  Arbitrary
translations stabilize $E_\tau$, $E_\sigma$, and $K$ semilinearly.
Therefore \eqref{eq:defect-cocycle} shows that $D_p$ is differentially
algebraic over $K$ for every $p\in P$.  If $h-h'\in P$, the same identity
shows that $D_h-D_{h'}$ is differentially algebraic over $K$.  Hence the
left-hand differential rank in \eqref{eq:coset-rank} is at most $\nu$.

It remains to prove independence for representatives
$h_1,\ldots,h_\nu$ of distinct nonzero classes.  Commuting translations in
\eqref{eq:psi-equations} gives
\begin{equation}
 \Deltaop_\sigma D_{h_i}=\Deltaop_{h_i}\lambda_\tau=:b_{\tau,i},
 \qquad
 \Deltaop_\tau D_{h_i}=\Deltaop_{h_i}\lambda_\sigma=:b_{\sigma,i}.
 \label{eq:paired-system}
\end{equation}
The right sides lie in $E_\tau$ and $E_\sigma$.  Thus
$F=K\langle D_{h_1},\ldots,D_{h_\nu}\rangle_\partial$ is stable under both
translations.  Every element of $F$ is $1$-periodic.  An element fixed by
$\T_\sigma$ and $\T_\tau$ has all of $P$ as periods, so
Lemma~\ref{lem:nondiscrete} gives
\[
 F^{\T_\sigma,\T_\tau}=\C.
\]

If the $D_{h_i}$ were differentially dependent over $K$,
Theorem~\ref{lem:ostrowski} applied to \eqref{eq:paired-system} would give
operators $L_i\in\C[\partial]$, not all zero, and $g\in K$ with
\begin{equation}
 (\T_\sigma-1)g=\sum_iL_i(b_{\tau,i}),\qquad
 (\T_\tau-1)g=\sum_iL_i(b_{\sigma,i}).
 \label{eq:paired-coboundary}
\end{equation}
Regard $g$ as a rational function on the product curve from
Lemma~\ref{lem:elliptic-fields}.  In the first equation, $\T_\sigma$ acts
only on the $\tau$-elliptic factor.  Lemma~\ref{lem:specialization} produces
$g_\tau\in E_\tau$ such that
\begin{equation}
 (\T_\sigma-1)g_\tau=\sum_iL_i(b_{\tau,i}).
 \label{eq:specialized}
\end{equation}

Choose $i$ with $L_i\ne0$.  On $\C/\Lambda_\tau$, $b_{\tau,i}$ has simple
poles exactly at $0$ and $-h_i$.  The $\T_\sigma$-orbit of $-h_i$ contains
neither $0$ nor any $-h_j$ with $j\ne i$: either collision would say that
$h_i\in P$ or $h_i-h_j\in P$.  Lemma~\ref{lem:highest-pole} therefore says
that the right side of \eqref{eq:specialized} has exactly one pole on that
orbit.  Translation by $\sigma$ is nontorsion on
$\C/\Lambda_\tau$, since otherwise $1,\sigma,\tau$ would be rationally
dependent.  Lemma~\ref{lem:pole} now contradicts \eqref{eq:specialized}.
Thus the representatives are differentially independent, proving the first
equality and rank value over $K$.

For the base $\C(z)$, both coefficient fields embed into $K(z)$.
Lemma~\ref{lem:elliptic-fields} and
Lemma~\ref{lem:da-preliminaries}(iii), applied once to $K\subset K(z)$ and
once to $\C(z)\subset K(z)$, show that the differential rank of every
$D$-packet is unchanged.  This proves the first equality for both bases.

Finally, logarithmic differentiation gives
\begin{equation}
 \frac{R_h'}{R_h}=D_h.
 \label{eq:quotient-logderiv}
\end{equation}
For either base $B$, the differential field generated by the $R_{h_i}$
contains all $D_{h_i}$, while each $R_{h_i}$ satisfies
$Y'-D_{h_i}Y=0$ over the field generated by the $D$-packet.  The intervening
extension is differentially algebraic, so the two packets have equal
differential transcendence degree.  This completes \eqref{eq:coset-rank}.
\end{proof}

\begin{corollary}[Exact spectrum and natural-base hypertranscendence]
\label{thm:elliptic-seed}
For $B=K$ or $B=\C(z)$,
\[
 D_h\text{ is differentially algebraic over }B
 \Longleftrightarrow
 R_h\text{ is differentially algebraic over }B
 \Longleftrightarrow h\in P.
\]
Moreover $\Psi$ and $\Gamma$ are hypertranscendental over $K$, and hence
also over $\C(z)$.
\end{corollary}

\begin{proof}
The equivalences are the one-element case of
Theorem~\ref{thm:coset-matroid}.  Choose $h\notin P$.  If $\Psi$ were
differentially algebraic over $K$, semilinear translation invariance and
field closure would make $D_h=\T_h\Psi-\Psi$ differentially algebraic, a
contradiction.  If $\Gamma$ were differentially algebraic over $K$, then
$\Psi=\Gamma'/\Gamma$ would be so as well.  If either function were
differentially algebraic over $\C(z)$, it would be differentially algebraic
over $K(z)$ and then, because $K(z)/K$ is differentially algebraic, over
$K$.  This gives the final claim.
\end{proof}

\section{Hypertranscendence climbs the face lattice}

For every nonempty subset $S\subseteq I$, let $G_S$ and $\Psi_S$ denote
the multiple elliptic gamma function and its logarithmic derivative for
the restricted period family $(\tau_i)_{i\in S}$.  Nishizawa's product
recursion is the exact identity
\begin{equation}
 G_S(z+\tau_q)=G_{S\setminus\{q\}}(z)\,G_S(z)
 \qquad(q\in S,\ |S|\geq2).
 \label{eq:face-product}
\end{equation}
Logarithmic differentiation gives
\begin{equation}
 \Deltaop_{\tau_q}\Psi_S=\Psi_{S\setminus\{q\}}.
 \label{eq:face-log}
\end{equation}

\begin{lemma}[Face escape]
\label{lem:escape}
If $q\notin S$ and $\Psi_S$ is hypertranscendental over $\C(z)$, then
$\Psi_{S\cup\{q\}}$ is hypertranscendental over $\C(z)$.
\end{lemma}

\begin{proof}
If $\Psi_{S\cup\{q\}}$ were differentially algebraic, translation by
$\tau_q$ and subtraction would preserve differential algebraicity.
Equation \eqref{eq:face-log} would then make $\Psi_S$ differentially
algebraic, a contradiction.
\end{proof}

\begin{lemma}[Rank-three minor extraction]
\label{lem:minor}
For a finite nonempty period tuple, $d_{\Q}(\tau)\geq3$ if and only if
there are distinct $q,r\in I$ such that
$1,\tau_q,\tau_r$ are rationally linearly independent.
\end{lemma}

\begin{proof}
One implication is immediate.  Conversely, extend the independent set
$\{1\}$ to a basis selected from the finite spanning family
$\{1\}\cup\{\tau_i:i\in I\}$.  Dimension at least three supplies two
period entries, with distinct labels because the selected vectors are
independent.
\end{proof}

\begin{proposition}[Rank-three-minor propagation]
\label{prop:propagation}
If $d_{\Q}(\tau)\geq3$, then $\Psi_I$ and $G_I$ are
hypertranscendental over $\C(z)$.
\end{proposition}

\begin{proof}
Choose $q,r$ as in Lemma~\ref{lem:minor}.  The two-period face
$\Psi_{\{q,r\}}$ is hypertranscendental by
Theorem~\ref{thm:elliptic-seed}.  List the labels in
$I\setminus\{q,r\}$ and adjoin them one at a time.  Lemma~\ref{lem:escape}
propagates hypertranscendence at each step, ending at $\Psi_I$.  If $G_I$
were differentially algebraic, then
$\partial G_I/G_I=\Psi_I$ would be differentially algebraic, a
contradiction.
\end{proof}

\begin{proof}[Proof of Theorem~\ref{thm:main}]
If $d_{\Q}(\tau)=2$, Proposition~\ref{prop:rank2} gives differential
algebraicity of both functions.  If $d_{\Q}(\tau)\geq3$,
Proposition~\ref{prop:propagation} gives hypertranscendence of both.
Since $d_{\Q}(\tau)\geq2$, the cases are exhaustive and disjoint.
\end{proof}

\section{Interpretation and relation to prior work}

The threshold in Theorem~\ref{thm:main} is forced by period geometry.
At augmented rank two, there is only one nonreal rational direction.  A
finite cover turns all nomes into powers of one nome, placing every factor
inside Kato's differential-algebraic phase.  At augmented rank at least
three, there exist distinct labels $q,r$ such that
$1,\tau_q,\tau_r$ are $\Q$-linearly independent.  Their two-period face has
the exact period-coset matroid of Theorem~\ref{thm:coset-matroid}.  Higher
periods cannot hide that obstruction because translation differences recover
lower faces exactly.  The proof flow is
\[
\begin{array}{ccccc}
d_{\Q}=2&\longrightarrow&\text{finite cover}&\longrightarrow&\text{Kato},\\[2mm]
d_{\Q}\geq3&\longrightarrow&\text{rank-three face}
&\longrightarrow&\text{joint descent}\longrightarrow
\text{orbit obstruction}\longrightarrow\text{face propagation}.
\end{array}
\]

Nishizawa introduced the hierarchy and its recursion
\cite{Nishizawa2001}.  Felder and Varchenko developed the automorphic
structure of the elliptic gamma function \cite{FelderVarchenko2000} and its
multiplication formulas \cite{FelderVarchenkoMultiplication}; Narukawa
established modular and integral representations for the multiple hierarchy
\cite{Narukawa2004}.  Kato proved hypertranscendence for multiple sine
functions with a nonreal period ratio \cite{Kato2019}, then differential
algebraicity of multiple elliptic gamma functions on the rational-ray locus
\cite[Theorem~1.1]{Kato2021}.  His Remark~1.2 explicitly left the
non-rational elliptic case open and anticipated exceptions.  Theorem
\ref{thm:main} closes that problem: the exceptions are exactly the
augmented-rank-two, equivalently nome-rank-one, tuples.

Parameterized difference Galois theory supplies the one-operator
differential-dependence criterion
\cite[Proposition~3.1]{HardouinSinger2008}.  Elliptic orbital residues were
developed in \cite[Appendix~B]{DreyfusHardouinRoquesSinger2018}, and
Arreche--Babbitt completed the summability obstruction by panorbital residues
\cite{ArrecheBabbitt2025}.  Theorem~\ref{lem:ostrowski} is best viewed as a
finite-family, common-trace formulation of the Ostrowski mechanism, not as a
replacement for that theory.  Lemmas~\ref{lem:pole} and
\ref{lem:highest-pole} use the known orbital mechanism in the special
one-private-orbit situation where no full summability algorithm is needed.
What is new in Theorem~\ref{thm:coset-matroid} is the simultaneous use of two
translations on the two elliptic factors to determine the complete pointed
differential matroid, over both the natural elliptic compositum and
$\C(z)$.

Complete linear-equation classifications over rational-function fields and
for global meromorphic shift equations appear in
\cite{AdamczewskiDreyfusHardouin2021,Dreyfus2026}.  Arreche, Dreyfus, and
Roques treated second-order equations and elliptic hypergeometric functions
\cite{ArrecheDreyfusRoques2021}.  De Shalit studied algebraic independence
for independent elliptic isogenies \cite{deShalit2025}, and de Shalit,
Hardouin, and Roques classified differentially algebraic solutions for a
single elliptic-curve isogeny \cite{deShalitHardouinRoques2024}.  These
results do not give either the all-order Nishizawa classification or the
two-factor translation-coset matroid above.

Recent geometric-family work \cite{MorainI,MorainII} studies modular
cocycles, partial modular symbols, and arithmetic applications.  Cone
generalizations \cite{Winding2016} alter the indexing geometry.  Neither
direction supplies a differential-algebraic classification of the standard
finite Nishizawa period tuple.  Theorem~\ref{thm:main} closes precisely that
classification and identifies Kato's anticipated exceptional
non-rational cases as the augmented-rank-two locus.  The classification is
qualitative for arbitrary complex input; Remark~\ref{rem:effective} records
its exact constructive reduction and why an algorithmic claim requires an
effective coefficient model.

\section{Conclusion}

For every nonempty finite upper-half-plane period tuple of the standard
Nishizawa hierarchy,
\[
 \boxed{
 \begin{array}{ccl}
 d_{\Q}(\tau)=2 &\Longleftrightarrow&
   \text{differential algebraicity},\\[2mm]
 d_{\Q}(\tau)\geq3 &\Longleftrightarrow&
   \text{hypertranscendence}.
 \end{array}}
\]
Equivalently, nome rank one is exactly the differential-algebraic phase and
nome rank at least two is exactly the hypertranscendental phase.  On the
rank-three seed, the stronger law
\[
 \dtrdeg_B B\langle D_{h_i}:1\leq i\leq m\rangle_\partial
 =\dtrdeg_B B\langle R_{h_i}:1\leq i\leq m\rangle_\partial
 =\#\{(h_i+P)\ne0\}
\]
holds for $B=E_\tau E_\sigma$ and $B=\C(z)$, with the set on the right
counted without repetition.  Thus the finite hierarchy has no unclassified
period configurations, and the seed has no unclassified finite translation
dependencies.

\section*{Declaration of generative AI and AI-assisted technologies in the
manuscript preparation process}

During preparation of this work, the author used GPT 5.6 Sol, Fable 5, and
Grok 4.5 for literature triage, drafting, and proof auditing.  The author
reviewed and edited the resulting content, checked the cited sources and
mathematical arguments, and takes full responsibility for the content of the
article.  These systems are not listed as authors because authorship requires
responsibility for the complete argument.

\section*{Data availability}

No empirical data or numerical computation was used in this work.

\begin{thebibliography}{99}

\bibitem{AdamczewskiDreyfusHardouin2021}
B. Adamczewski, T. Dreyfus, and C. Hardouin,
\emph{Hypertranscendence and linear difference equations},
J. Amer. Math. Soc. \textbf{34} (2021), 475--503,
\url{https://doi.org/10.1090/jams/960}.

\bibitem{ArrecheBabbitt2025}
C. E. Arreche and M. W. Babbitt,
\emph{Panorbital residues and elliptic summability},
arXiv:2508.18247, 2025,
\url{https://doi.org/10.48550/arXiv.2508.18247}.

\bibitem{ArrecheDreyfusRoques2021}
C. E. Arreche, T. Dreyfus, and J. Roques,
\emph{Differential transcendence criteria for second-order linear difference
equations and elliptic hypergeometric functions},
J. \`Ec. polytech. Math. \textbf{8} (2021), 147--168,
\url{https://doi.org/10.5802/jep.143}.

\bibitem{deShalit2025}
E. de Shalit,
\emph{Algebraic independence and difference equations over elliptic function
fields},
Ann. Inst. Fourier \textbf{75} (2025), 1509--1554,
\url{https://doi.org/10.5802/aif.3694}.

\bibitem{deShalitHardouinRoques2024}
E. de Shalit, C. Hardouin, and J. Roques,
\emph{Hypertranscendence and $q$-difference equations over elliptic function
fields},
arXiv preprint arXiv:2409.10092, 2024,
\url{https://doi.org/10.48550/arXiv.2409.10092}.

\bibitem{Dreyfus2026}
T. Dreyfus,
\emph{Hypertranscendence and linear difference equations, the exponential
case},
J. Algebra \textbf{686} (2026), 775--792,
\url{https://doi.org/10.1016/j.jalgebra.2025.08.025}.

\bibitem{DreyfusHardouinRoquesSinger2018}
T. Dreyfus, C. Hardouin, J. Roques, and M. F. Singer,
\emph{On the nature of the generating series of walks in the quarter plane},
Invent. Math. \textbf{213} (2018), 139--203,
\url{https://doi.org/10.1007/s00222-018-0787-z}.

\bibitem{FelderVarchenko2000}
G. Felder and A. Varchenko,
\emph{The elliptic gamma function and
$\mathrm{SL}(3,\Z)\ltimes\Z^3$},
Adv. Math. \textbf{156} (2000), 44--76.

\bibitem{FelderVarchenkoMultiplication}
G. Felder and A. Varchenko,
\emph{Multiplication formulae for the elliptic gamma function},
in \emph{Noncommutative Geometry and Representation Theory in Mathematical
Physics}, Contemp. Math. \textbf{391}, Amer. Math. Soc., 2005, 69--73;
\url{https://doi.org/10.48550/arXiv.math/0212155}.

\bibitem{HardouinSinger2008}
C. Hardouin and M. F. Singer,
\emph{Differential Galois theory of linear difference equations},
Math. Ann. \textbf{342} (2008), 333--377,
\url{https://doi.org/10.48550/arXiv.0801.1493}.

\bibitem{Kato2021}
M. Kato,
\emph{Differential algebraicity of the multiple elliptic gamma function
for a rational period},
Funkcial. Ekvac. \textbf{64} (2021), 225--235,
\url{https://doi.org/10.1619/fesi.64.225}.

\bibitem{Kato2019}
M. Kato,
\emph{Hypertranscendence of the multiple sine function for a complex period},
Proc. Japan Acad. Ser. A Math. Sci. \textbf{95} (2019), 16--19,
\url{https://doi.org/10.3792/pjaa.95.16}.

\bibitem{MorainI}
P. L. L. Morain,
\emph{Geometric families of multiple elliptic Gamma functions and
arithmetic applications, I},
arXiv:2510.16515, 2025.

\bibitem{MorainII}
P. L. L. Morain,
\emph{Geometric families of multiple elliptic Gamma functions and
arithmetic applications, II},
arXiv:2602.06561, 2026.

\bibitem{Narukawa2004}
A. Narukawa,
\emph{The modular properties and the integral representations of the
multiple elliptic gamma functions},
Adv. Math. \textbf{189} (2004), 247--267.

\bibitem{Nishizawa2001}
M. Nishizawa,
\emph{An elliptic analogue of the multiple gamma function},
J. Phys. A \textbf{34} (2001), 7411--7421.

\bibitem{Winding2016}
J. Winding,
\emph{Multiple elliptic gamma functions associated to cones},
Adv. Math. \textbf{325} (2018), 56--86.

\end{thebibliography}

\end{document}
